\documentclass[../../../include/open-logic-section]{subfiles}

\begin{document}

\olfileid{sth}{choice}{banach}
\olsection{باناخ–تارسکی دا پیراڈوکس}

جیویں بُولوس نے بدل رکھن دے اصول دا \emph{باہرلا} جواز دین دی کوشش کیتی، اسیں چُون دا جواز وی اویویں باہرلے پہلوآں توں دین دی کوشش کر سکدے آں (\olref[replacement][extrinsic]{sec} ویکھو)۔ آخر چُون نوں منّن دے کئی من پسند نتیجے نیں۔ ہر کارڈینل عدد دا موازنہ، ہر سیٹ نوں حُسنِ ترتیب، کارڈینل عدداں نوں خاص قسم دے ترتیبی عدد سمجھ سکنا، تے ہور وی۔

پر کدی کدی ایہ دعویٰ ہُندا اے کہ چُون دے \emph{ناپسند} نتیجے وی نیں۔ ایہ دعویٰ بہت کرکے \cite{BanachTarski1924} دے اک نتیجے کارن ہُندا اے۔

\begin{thm}[باناخ–تارسکی دا پیراڈوکس ($\ZFC$ وچ)]
کسے وی ٹھوس گولے نوں محدود گنتی دے ٹکڑیاں وچ ونڈیا جا سکدا اے۔ اوہ ٹکڑے فیر (گھما کے تے تھاں بدل کے) جوڑے جا سکدے نیں تاں جو اوس گولے دیاں دو نقلاں بن جان۔
\end{thm}
\noindent
پہلی نظر وچ ایہ گل حیرانی والی اے۔ صاف اے کہ دو گولیاں دا حجم پہلے گولے دے حجم توں \emph{دوگنا} اے۔ پر شکل قائم رکھن والیاں حرکتاں، گھماؤ تے تھاں بدلنا، حجم نہیں بدلدیاں۔ سو ایویں لگدا اے جیویں باناخ–تارسکی سانوں جادو نال نواں مادہ پیدا کرن دیندا اے۔

گل ہور وی حیرانی والی اے۔\footnote{\citet[Theorem 3.12]{Wagon2016} ویکھو۔} ایہو جیہی دلیل وکھاندی اے کہ مٹر دے اک دانے نوں محدود گنتی دے ٹکڑیاں وچ ونڈ کے اوہناں نوں فیر (گھما کے تے تھاں بدل کے) جوڑیا جا سکدا اے، تاں جو بِگ بین دی شکل تے وڈائی والی شے بن جائے۔

پر ایہناں وچوں کوئی گل اکلّے $\ZF$ وچ نہیں چل دی۔\footnote{بھاویں باناخ–تارسکی اوہناں اصولاں نال ثابت ہو سکدا اے جیہڑے چُون توں ٹھیک ٹھیک کمزور نیں۔ \citet[303]{Wagon2016} ویکھو۔} سو ساڈے اگے اک فیصلہ اے۔ چُون ردّ کریے یا ایس «پیراڈوکس» نوں قبول کرنا سکھیے۔

اسیں آکھاں گے کہ سانوں «پیراڈوکس» قبول کرنا سکھنا چاہیدا اے۔ اصل وچ سانوں ایہ کوئی خاص پیراڈوکس لگدا ای نہیں۔ خاص طور اُتے سانوں کوئی وجہ نہیں دِسدی کہ ایہ ہِیٹھلے نتیجیاں نالوں ودھ یا گھٹ پیراڈوکس کیوں اے۔\footnote{\citet[276--7]{Potter2004}، \citet[16]{Weston2003}، تے \citet[31, 308--9]{Wagon2016} ہور مثالاں ورت کے ایہو جیہیاں گلاں کردے نیں۔}
\begin{enumerate}
	\item وقفے $(0,1)$ وچ اینے ای نقطے نیں جِنّے $\Real$ وچ نیں۔
	\\\emph{ثبوت}۔ $\tan(\pi(r-\nicefrac{1}{2})))$ ویکھو۔
	\item اک لکیر وچ اینے ای نقطے نیں جِنّے اک مربع وچ۔
	\\See \olref[his][set][pathology]{sec} تے \olref[his][set][cantorplane]{sec} ویکھو۔
	\item تھاں بھرن والیاں منحنی لکیراں موجود نیں۔
	\\See \olref[his][set][pathology]{sec} تے \olref[his][set][hilbertcurve]{sec} ویکھو۔
\end{enumerate}
ایہناں تِنّاں نتیجیاں وچوں کسے لئی چُون دی لوڑ نہیں۔ اصل وچ ہُن اسیں بس ایہناں نوں ریاضی دیاں حیرانی والیاں، سوہنیاں گلاں سمجھدے آں۔ شاید باناخ–تارسکی دے پیراڈوکس نال وی سانوں ایہو رویہ رکھنا چاہیدا اے۔

ایتھے اک فنی گل دا دھیان دینا ضروری اے، پر ایہ بس ٹھنڈے دل نال سوچن دی گل اے۔ شکل قائم رکھن والیاں حرکتاں حجم وی قائم رکھدیاں نیں۔ ایس لئی پنجے\footnote{اسیں پیراڈوکس نوں «محدود گنتی دے ٹکڑیاں» نال بیان کیتا سی۔ اصل وچ \citet{Robinson1947} نے ثابت کیتا کہ گولے نوں \emph{پنج} ٹکڑیاں وچ ونڈنا کافی اے (ایہتوں گھٹ وچ نہیں ہو سکدا)۔ ثبوت لئی \citet[pp.~66--7]{Wagon2016} ویکھو۔} ٹکڑے، جنہاں وچ گولا ونڈیا اے، سارے دے سارے \emph{ناپن جوگ} نہیں ہو سکدے۔ موٹی گل ایہ اے کہ ایہناں وکھ وکھ ٹکڑیاں دا حجم مقرر کرن دا کوئی مطلب نہیں۔ ایہناں نوں نقطیاں دے ایہو جیہے «لامتناہی کھلارے» سمجھو جنہاں دی تصویر ذہن وچ نہیں بن سکدی۔ ہُن ایہ ہو سکدا اے کہ ایہو جیہے «لامتناہی کھلرے» سیٹ سوچنا اوپرا لگے۔ پر اوہناں دی موجودگی \stagesacc{} وچ دتے حکم توں نکلدی لگدی اے کہ پہلاں موجود سیٹاں دے \emph{سارے ممکن} اکٹھ بناؤ۔

جے ایہناں وچوں کوئی گل تہانوں نہیں مناؤندی، تاں باناخ–تارسکی دا پیراڈوکس قبول کرن دے حق وچ اک آخری (باہرلی) دلیل ویکھو۔ ایہ توں ریاضی دا اج تک دا سبھ توں وڈیا لطیفہ جھٹ نکلدا اے:
\begin{enumerate}
	\item[] \emph{سوال}۔ «باناخ–تارسکی» دے اکھر اگے پچھے کر کے کیہ بن دا اے؟
	\item[] \emph{جواب}۔ «باناخ–تارسکی باناخ–تارسکی»۔
\end{enumerate}

\end{document}
